\documentclass[11pt]{article}

\title{InjectRAG: Indirect Prompt Injection and Spotlighting Defenses}
\author{Your Name}
\date{\today}

\begin{document}
\maketitle

\begin{abstract}
% Summarize the objective, attack setting, evaluated defenses, and measured outcome.
\end{abstract}

\section{Introduction}
\subsection{Problem Statement}
\subsection{Project Objectives}

\section{Background}
\subsection{Retrieval-Augmented Generation}
\subsection{Indirect Prompt Injection}
\subsection{Spotlighting as a Defense}

\section{System Design}
\subsection{RAG Architecture and Data Flow}
\subsection{Corpus and Document Roles}
\subsection{Attacker, Victim, and System Interfaces}
% Add an architecture diagram showing ingestion, retrieval, generation, and visible outputs.

\section{Threat Model and Attack Design}
\subsection{Attacker Capabilities and Assumptions}
\subsection{Victim Scenario and Attacker Objective}
\subsection{Payload Structure}
\subsection{Attack Procedure}
% Use numbered steps; cite screenshots of the submitted content and victim screen.
\subsection{Expected Outcome}

\section{Defense Design}
\subsection{Delimiting}
\subsection{Datamarked Spotlighting}
\subsection{Expected Defense Behavior}
% Explain the prompt or document transformation for each defense.

\section{Experimental Setup}
\subsection{Hardware and Software Environment}
\subsection{Corpus, Document, and Query Counts}
% Distinguish source documents from retrieved chunks and report counts per split.
\subsection{Models and Parameters}
% Report embedding model, generator, retrieval settings, and generation settings.
\subsection{Prompt Templates and Versions}
% Summarize prompts here; put complete text in the appendix.
\subsection{Run Conditions and Evaluation Metrics}
% Define success criteria and denominators before presenting results.

\section{Results and Observations}
\subsection{Clean RAG Baseline}
\subsection{Attack Without Defense}
\subsection{Attack With Delimiting}
\subsection{Attack With Datamarked Spotlighting}
\subsection{Comparison Across Conditions}
% One table: condition, attempts, successes, success rate, and observed behavior.
\subsection{Outputs on Attacker, Victim, and System Sides}
% Include representative screens or logs, with captions and run identifiers.

\section{Case Studies}
\subsection{Successful Attack}
\subsection{Failed Attack}
\subsection{Defense Success or Partial Failure}
% For each case: payload, expected result, retrieved evidence, observed outputs,
% actual result, and explanation. Include only cases supported by recorded runs.

\section{Analysis}
\subsection{Why Attacks Succeeded or Failed}
\subsection{Expected Versus Actual Outcomes}
\subsection{Effectiveness of Delimiting and Datamarking}
\subsection{Assumptions and Limitations}

\section{Conclusion}
\subsection{Summary of Findings}
\subsection{Lessons Learned and Future Work}

\appendix
\section{Full Prompt Templates}
\section{Attack Payloads}
\section{Additional Screenshots and Logs}
\section{Raw Result Tables}

\end{document}
